% SPDX-License-Identifier: CC-BY-SA-4.0
%
% Course-companion paper: teaching modern OS design through
% kernel-extension VMs.
%
% Self-contained workshop / educational-track paper; not tied to a
% specific venue's class file.  Use article.cls so it compiles on a
% stock TeX Live install.
\documentclass[11pt,a4paper]{article}

\usepackage[T1]{fontenc}
\usepackage[utf8]{inputenc}
\usepackage{lmodern}
\usepackage{microtype}
\usepackage{amsmath}
\usepackage{listings}
\usepackage{xcolor}
\usepackage{hyperref}
\usepackage{cite}
\usepackage[a4paper,margin=2.4cm]{geometry}
\usepackage{enumitem}

\hypersetup{
  colorlinks=true,
  linkcolor=blue!50!black,
  urlcolor=blue!50!black,
  citecolor=blue!50!black,
}

\lstdefinelanguage{rvasm}{
  morekeywords={addi,add,sub,lw,sw,lbu,lhu,beq,bne,blt,bge,jal,jalr,
                lui,auipc,ecall,fence},
  sensitive=true,
  morecomment=[l]{;},
  morecomment=[l]{//},
}

\lstset{
  basicstyle=\ttfamily\footnotesize,
  keywordstyle=\bfseries,
  commentstyle=\color{gray},
  numbers=left,
  numberstyle=\tiny\color{gray},
  numbersep=8pt,
  frame=single,
  framerule=0.3pt,
  xleftmargin=1.5em,
}

\title{Teaching Modern Operating Systems Design Through Kernel-Extension VMs:\\
       A MERLIN-V Course Companion}
\author{Weqaar Janjua\\
        \texttt{weqaar.janjua@packetfive.com}\\
        PacketFive --- University of Limerick}
\date{2026}

\begin{document}
\maketitle

\begin{abstract}
We argue that the kernel-extension virtual machine ---
exemplified by eBPF in Linux --- is the single best concrete
vehicle for teaching modern operating-systems design at the
graduate level.  It exercises, in one project, every load-bearing
concept the canonical OS curriculum covers in isolation: machine
representation, abstract interpretation, JIT codegen, virtual
memory and W$\oplus$X, RCU, I-cache coherence, syscall ABI design,
sandboxing without an MMU, cross-OS portability, and the
discipline of designing an interface that cannot be broken in a
future release.
We present \textbf{MERLIN-V}, an in-kernel JIT VM whose bytecode
is a restricted profile of the RISC-V ISA, and the course we
built around it.  The course delivers eleven labs that take a
student from RV32 hand-encoding to a working Linux out-of-tree
kernel module that loads, verifies, and JIT-executes the same
\texttt{.merlin.o} that also runs on a 400~KiB SRAM RISC-V MCU
under Zephyr.  Every lab is autograded by a CI pipeline that
overlays canonical instructor files on top of the student's
fork before grading, preventing test-suite tampering.  Our claim
is that this is not a niche project: it is the cleanest pedagogy
for the next generation of OS curricula, because the
\emph{interesting} OS-level engineering happening in industry
today is almost entirely about extension and sandboxing,
not about classical scheduler design.
\end{abstract}

% ---------------------------------------------------------------
\section{Motivation}
\label{sec:motivation}

The canonical operating-systems curriculum covers the process,
virtual memory, the file system, the scheduler.  This is the
material in OSTEP~\cite{ostep} and in every undergraduate OS
textbook.  It is correct material, but it covers ground that the
industry has solved at production scale by the early 2000s, and
the open research and engineering questions in OS today look
different.

The frontier where senior kernel engineers actually work is
\emph{controlled extension}: how do we let third-party code
modify or observe a running production kernel safely?  eBPF is
the answer Linux landed on.  ChromeOS sandboxing, Cilium's
networking, Meta's tracing infrastructure, and a large fraction
of modern observability tooling depend on eBPF programs the user
\emph{loads at runtime}, the kernel \emph{verifies} as safe, and
the JIT compiles to host machine code.

Yet eBPF rarely appears in OS curricula.  When it does, it
appears at the level of ``here is \texttt{bpftrace}'', not at the
level of ``here is how a verifier proves a program terminates.''
The lab-companion project of an OS course typically remains a
toy scheduler or a toy filesystem.

We argue this is a missed opportunity.  A modern kernel-extension
VM is, simultaneously:
\begin{itemize}[leftmargin=*]
  \item \textbf{A compiler problem.}  The bytecode must be
    chosen, the verifier built, and the JIT emit machine code.
  \item \textbf{A program-analysis problem.}  Soundness is
    non-negotiable.  Abstract interpretation, type-state, and
    SMT-adjacent reasoning are required.
  \item \textbf{A kernel-internals problem.}  Loading,
    refcounting, RCU, anon-inode fds, I-cache coherence on
    multi-core systems, and W$\oplus$X all show up.
  \item \textbf{A systems-design problem.}  The UAPI must
    survive Linux's ``never break userspace'' contract.  The
    helper-ABI must be efficient, verifiable, and stable.
  \item \textbf{A cross-OS portability problem.}  The same
    bytecode should be loadable into an embedded RTOS on a
    microcontroller.
\end{itemize}

No other single project hits all five.

% ---------------------------------------------------------------
\section{MERLIN-V in One Page}
\label{sec:merlin}

MERLIN-V is a clean-room in-kernel JIT VM whose bytecode is a
restricted profile of the RISC-V ISA.  The premise is small:
eBPF's bytecode is a custom 64-bit RISC-like ISA that exists
because, in 2014, there was no widely available, free, simple,
64-bit RISC ISA with production toolchains.  That precondition
no longer holds.  RISC-V~\cite{riscv-isa-manual} is ratified,
ABI-stable, multi-vendor, and supported in GCC, LLVM, glibc,
musl, Linux, and Zephyr~\cite{zephyr-rtos}.

The consequence is structural, not incremental:
\begin{itemize}[leftmargin=*]
  \item On a RISC-V host the JIT is a pass-through step --- the
    verified image \emph{is} native machine code.  No
    translation.
  \item On a non-RISC-V host (x86\_64, arm64, \ldots) the JIT
    still exists but its input is standard RISC-V machine code
    rather than a bespoke bytecode.
  \item The toolchain story collapses to ``stock RISC-V GCC
    plus section macros.''  No bespoke LLVM-BPF backend; no
    parallel \texttt{gcc-bpf}.
\end{itemize}

The full design record is at~\cite{merlinv-design}.  For the
purposes of this paper, what matters is that MERLIN-V exposes
every interesting OS-level problem mentioned in
\S\ref{sec:motivation}, without the historical baggage of eBPF's
custom ISA.  That makes it pedagogically tractable.

% ---------------------------------------------------------------
\section{A Worked Example}
\label{sec:example}

Throughout the course, students return to a single seven-instruction
program.  In RV32 assembly it is:

\begin{lstlisting}[language=rvasm,caption={Packet-EtherType classifier.},label=lst:classifier]
classify:
    lbu  a1, 12(a0)       ; a1 = pkt[12]
    addi a2, x0, 8        ; a2 = 0x08 (IPv4 high byte)
    beq  a1, a2, .pass
    addi a0, x0, 1        ; return 1 (DROP)
    jalr x0, ra, 0
.pass:
    addi a0, x0, 2        ; return 2 (PASS)
    jalr x0, ra, 0
\end{lstlisting}

Listing~\ref{lst:classifier} is the smallest non-trivial
verifier exercise.  The verifier must reason that
\texttt{a0} on entry is a typed pointer (the program context);
that \texttt{lbu} loading through it is therefore safe; that
\texttt{a2} on the \texttt{beq} is a compile-time-known constant
(so the verifier can fold the branch in the abstract domain);
that the back-edge-free CFG terminates; and that neither
\texttt{ecall} nor any privileged instruction appears.

The same 28~text bytes assemble into a complete
\texttt{.merlin.o} (an ET\_REL Elf32 for the rtos-rv32 profile
running on Zephyr/ESP32-C3-DevKitM-1, and an ET\_REL Elf64 for
the linux-rv64 profile running on the Linux kernel module on
the MPFS Icicle Kit).  This is the cross-OS portability
demonstration the course closes with.

% ---------------------------------------------------------------
\section{Course Structure}
\label{sec:course}

The course is organised into six modules covering eleven labs.
Modules build on each other; there is no calendar pacing in the
syllabus because the same material works for a quarter, a
semester, or a graduate-level block course.

\begin{enumerate}[leftmargin=*]
  \item \textbf{Setup} --- environment, cross-toolchain, board
    bring-up (Lab 00).
  \item \textbf{Context} --- read an eBPF object end-to-end
    (Lab 01); learn RV32 assembly by hand (Lab 02).
  \item \textbf{User-space MERLIN-V core} --- build a user-space
    interpreter (Lab 03), a minimal verifier (Lab 04), and a
    section-validator / objtool (Lab 05).
  \item \textbf{JITs} --- pass-through (Lab 06) and host
    RV$\rightarrow$x86\_64 (Lab 07).
  \item \emph{Midterm: a written design exam, no code.}
  \item \textbf{In-kernel runtime} --- an out-of-tree Linux
    kernel module that loads, verifies, JITs, and runs a
    \texttt{.merlin.o} (Lab 08); the same runtime ported to
    Zephyr (Lab 09).
  \item \textbf{Hardware bring-up + capstone} --- run the
    course's worked example on real hardware (Lab 10), then
    a student-chosen capstone project (Lab 11).
\end{enumerate}

Lab~03 is the load-bearing one.  By the end of it, a student
who started knowing only C has written a verifier-friendly RV32
interpreter that runs Listing~\ref{lst:classifier}.  Lab~04
turns it into an abstract interpreter; the leap from concrete
to abstract semantics is the single highest-leverage moment in
the course.

% ---------------------------------------------------------------
\section{Autograder Pipeline}
\label{sec:autograder}

Each lab ships with a Makefile, a SKELETON / PARTIALLY-PROVIDED
implementation file, and a \texttt{tests/} directory.  The
autograder is a 90-line shell script that builds the student's
tree, runs \texttt{make test}, scrapes ``N/M tests passed,''
and emits one JSON record per lab:

\begin{lstlisting}[caption={Autograder JSON record.}]
{"lab":"lab-03","build":"ok","pass":7,"total":7,"score":100}
\end{lstlisting}

The same autograder runs locally and in CI.  In CI (GitHub
Actions and Azure DevOps pipelines both ship) the harness
overlays the canonical instructor files --- \texttt{Makefile},
\texttt{tests/}, and any header marked \texttt{PROVIDED} ---
on top of the student's push before running the grader.  This
eliminates an entire class of academic-integrity bugs: a
student cannot pass by replacing \texttt{tests/run-all.sh} with
\texttt{echo 100/100 tests passed; exit 0}, because the overlay
silently restores the canonical version first.

We verified this empirically.  A simulated cheating student
whose \texttt{run-all.sh} fakes a perfect score gets:
\begin{itemize}[leftmargin=*]
  \item Without overlay: \texttt{score=100} (cheat succeeds).
  \item With overlay (the CI path):  \texttt{score=14}
    (their actual TODO completion).
\end{itemize}

The pipeline posts the grade table as a Markdown comment on the
student's PR and uploads the JSON manifest as a workflow
artifact.  Instructors aggregate the artifacts into a
gradebook; no per-student manual grading is required for
Labs~02--09.  Labs~10 and~11 require hardware-in-the-loop
artifacts (boot logs, demo recordings) and are graded
manually.

% ---------------------------------------------------------------
\section{What This Project Teaches That a Toy Kernel Does Not}
\label{sec:contrast}

A toy-kernel project --- xv6, JOS, the OSTEP exercises --- gives
students experience with system primitives but stops short of
the engineering culture of modern kernel development:
\begin{itemize}[leftmargin=*]
  \item \textbf{Stable UAPI.}  Toy kernels are free to break
    their own ABI every assignment.  MERLIN-V's
    \texttt{union merlin\_attr} is designed to honour the same
    ``never break userspace'' rule eBPF lives under.  Students
    learn how to extend a structure additively, with
    zero-initialised tail regions and explicit version
    negotiation.
  \item \textbf{Verifier soundness.}  Students discover, by
    building it, that the verifier is the single trusted
    component.  A bug in the JIT corrupts one program; a bug
    in the verifier compromises the kernel.  The asymmetry is
    not abstract; their lab points it out concretely.
  \item \textbf{RCU and refcounting.}  The kernel module
    refcounts program objects through anon-inode fds; hot-swap
    of a program requires an RCU grace period.  Students
    encounter these as real synchronisation primitives, not
    textbook artefacts.
  \item \textbf{Cross-OS portability.}  The same
    \texttt{.merlin.o} runs on Linux and Zephyr.  Students see
    that the kernel's calling convention, page-table
    discipline, and I-cache coherence are decoupled from the
    \emph{language} the program is written in.
  \item \textbf{CI and academic integrity.}  The pipeline
    teaches students how production code is integrated --- pull
    requests, automated grading, locked-down test artefacts
    --- in a context where the gamesmanship has actual
    consequences (their grade).
\end{itemize}

% ---------------------------------------------------------------
\section{Related Work in Pedagogy}
\label{sec:related}

The closest comparable lab projects in OS pedagogy are MIT's
xv6, Stanford's Pintos, and various RISC-V CPU-in-Verilog
courses (e.g.\ Onur Mutlu's at ETH Zurich).  These are
excellent at what they do, but they all stop at the level of
``build a kernel.''  None of them touches kernel
\emph{extension}, which is where the engineering frontier sits.

The closest peer in the JIT-pedagogy literature is the work
emerging from the WebAssembly community~\cite{wasm-spec}, which
shares the ``verified bytecode + JIT'' design.  WebAssembly's
pedagogical material is, however, almost entirely
user-space-oriented.  The kernel context --- RCU, I-cache, IPI,
UAPI stability --- is missing.  MERLIN-V fills that gap.

% ---------------------------------------------------------------
\section{Open Questions}
\label{sec:open}

The course is in active use; we close with the questions we
expect future iterations to resolve.

\begin{itemize}[leftmargin=*]
  \item \textbf{Verifier sharing with \texttt{kernel/bpf/}.}
    Should the MERLIN-V verifier share its type-state plumbing
    with eBPF's?  The pedagogical answer is currently
    ``no'' --- duplication makes the abstract interpreter
    legible.  Production answer may differ.
  \item \textbf{The midterm.}  We treat the midterm as a
    written design exam (no code) because the early labs do
    not need integration testing.  Whether this is the right
    pedagogical choice for a remote / asynchronous version of
    the course is open.
  \item \textbf{Capstone open-endedness.}  Lab~11 is
    deliberately under-specified.  Past students have produced
    an arm64 JIT, a verifier extension for bounded loops, a
    \texttt{merlin-trace} kfunc registration framework, and a
    PolarFire fabric soft-core offload demo.  An equally good
    choice may be a tighter capstone with a single canonical
    deliverable.
\end{itemize}

% ---------------------------------------------------------------
\section{Conclusion}

The kernel-extension VM is the right pedagogical vehicle for
the next generation of OS curricula because it is what kernel
engineers actually build.  MERLIN-V removes a barrier eBPF
imposes --- the custom bytecode --- and brings the entire
toolchain story into a domain (RISC-V) that students will
encounter again in compilers, computer architecture, and
embedded courses.  Eleven labs, an autograded CI pipeline,
two reference platforms, and the same seven-instruction worked
example across all of them.

The course materials, lab skeletons, and CI pipelines are open
source at~\cite{merlinv-design} under CC-BY-SA-4.0 (text) and
GPL-2.0-only / Apache-2.0 (code).  We invite instructors at
other institutions to adopt, fork, and contribute back.

% ---------------------------------------------------------------
\bibliographystyle{plain}
\bibliography{refs}

\end{document}
